\section{Historical latent-class analysis}
\label{sec:lca}
This section documents the original exploratory analysis and its numerical reconstruction. The paper reports observed profiles and does not use historical classes in its outcome analysis. Section~\ref{sec:revised-lca} documents the separate analysis with revised indicators.
\subsection{The archived candidate}
The historical analysis compared four minimal-binary and four mixed-input specification families (M1--M4 and H1--H4), each at two to five classes. A separate preferred-final execution, based on the H1 four-class family, is not interchangeable with the original four-class grid attempt. It used six three-category indicators for the historical vulnerability, digital capability, formal financial access, formal-system use, service quality and digital remittance/payment-use constructs (IVS, ICDP, IAFF, IUOF, OQI and IURD), and a binary adverse-experience indicator (IEDF). The archived command used a separate ordered-logit measurement equation for each ternary item and a binary logit for IEDF, with parameters varying by class. The indicators were recoded from historical composites rather than collected as seven standalone survey items. The retained complete-case fit covered 423 records. The source sorted by the submission key, used the Stata version 16 declaration, set seed 6529004, called \texttt{randomid draws(600)}, and specified \texttt{difficult} and an iteration limit of 8,000. These are provenance facts, not settings selected for an academic refit.

The archived IEDF binary is a source-coded adverse-experience proxy, not a literal split above the exact median. The historical Stata macro stored a cutoff just below the exact stored-score median, placing 147 tied records on the upper side. The audit counted 269 positives under the implemented rule and 122 under a strict split at the exact median. The historical model used the first definition. In separate sensitivity fits on the full and 407-case matched samples, using the exact-median rule instead changed aligned four-class IEDF probabilities by as much as about 52 percentage points. The fits still failed the numerical admissibility checks. The difference is a change in scoring, not in reported experiences.

\subsection{What the model assumes}
Let $Y_i=(Y_{i1},\ldots,Y_{iJ})$ be the categorical indicator pattern and $C_i$ the latent class. With local independence, a $K$-class model represents the probability of an observed pattern as
\[
 \Pr(Y_i=y_i)=\sum_{c=1}^K \pi_c\prod_{j=1}^J p_{jc}(y_{ij}), 
 \qquad \sum_c\pi_c=1,
\]
where $p_{jc}(y)$ is the class-specific probability of response $y$ for item $j$. The observed-data log likelihood sums the logarithm of this pattern probability over the model sample. The class weight $\pi_c$ is a model mixing probability. A respondent's posterior membership weight, if computed, is $\pi_c\prod_j p_{jc}(y_{ij})/\Pr(Y_i=y_i)$. It is not an observed characteristic. Neither mixing weights nor sums of posterior probabilities should be confused with counts assigned to the highest-probability class.

For the preferred-final specification, the six ternary items and one binary item allow $3^6(2)=1{,}458$ theoretical patterns, not necessarily that many observed patterns. Four classes allow $(4-1)+4[6(3-1)+(2-1)]=55$ free probability coordinates, consistent with 55 free coordinates across 56 stored native slots. These counts do not establish identification. A simultaneous permutation of class labels leaves the distribution unchanged. Beyond that invariance, identification and reliable estimation depend on sufficiently distinct, informative response profiles and a defensible measurement model \citep{allmanMatiasRhodes2009}. Local independence can be implausible when historical indices share items, when skip patterns jointly determine what is observable, or when recoding implicitly measures common conditions. Classes can partly reflect that dependence, rather than stable types of people. Category order does not imply that a higher value is desirable across all the domains.

\subsection{What historical reproducibility establishes}
Among converged H1 grid fits, four classes had lower BIC than two; three and five classes did not converge, and the grid had no one-class baseline (Table~\ref{tab:historical-h1}). The two-class fit had slightly higher entropy. The archived screen admitted only four- and five-class fits as reporting candidates, so it cannot establish that four classes were optimal. The 600 random initial draws are part of one initialization strategy, not 600 independently optimized replications.

\input{"\CFIPaperDir supplement/historical-lca-enumeration.tex"}

The saved-state comparison matched 3,199 captured native fields per state, including the retained convergence status. The audit also found an item-category probability near the boundary (about $9.2\times 10^{-15}$), a native cutpoint with a standard error about $1.73$ million, and severe scaling problems in the raw covariance coordinates. Better-scaled probability coordinates had full local numerical rank at the retained tolerances. That finding neither proves global uniqueness nor guarantees regular Wald inference or stable assignment. A stored convergence flag and exact archived numerical equality answer a computational question, not whether the measurement model fits observed patterns adequately or should be reported as an academic finding. The class-profile audit below used the saved coefficients and frequency-counted response patterns. It did not refit the model or make new respondent-level assignments.

\input{"\CFIPaperDir supplement/historical-h1-profiles.tex"}

\subsection{Why the paper does not adopt the historical classes}
The historical H1 input has 291 observed patterns among the 1,458 possible; 220 patterns occur once. That sparsity calls for care with asymptotic goodness-of-fit tests. The pre-fit correlations and Cramér's~V summaries do not test local independence after fitting. In particular, the formal-access, financial-use and service-quality measures share survey routing, which could induce residual dependence. The historical profile displays report posterior-weighted or modal-assignment proportions in observed categories. Those quantities are distinct from native model-implied response probabilities, and none is independent validation when the item helped define the classes.
The first question is whether every class-defining item measures the intended construct under the adopted recoding and applicable denominator. Several historical composites have been split into separate items in the academic plan, including IURD and IEDF. Keeping a historical name does not authorize substituting a different construct into the same fit. The second question is whether candidate counts are compared within a common indicator family using the same likelihood and eligible records. Information criteria are useful within that comparison but do not select a preferred model across different response variables. Convergence, comparable maxima across starts, profile separability, plausible local dependence, uncertainty and substantive interpretability all matter; a rule based only on entropy, a smallest-class cutoff or the numerical maximum is insufficient \citep{nylundEtAl2007}. An ordinary chi-square likelihood-ratio reference is not assumed when an additional mixture component is at the boundary.

A downstream relationship using an item that helped define the classes is partly mechanical; it is not external validation. For a genuinely external outcome, treating an assigned class as observed ignores classification error. A one-step extension can change the profiles themselves; an error-adjusted multi-step approach needs assumptions suited to the outcome and design \citep{vermunt2010}. Neither approach addresses the overlap built into the historical measures. The observed profiles in the paper answer a narrower question directly: which combination of onboarding conditions did respondents report? This does not rule out latent heterogeneity; it avoids treating this archived solution as evidence for distinct types.
\paragraph{A historical predictor that is not independent of the class construction.}
The vulnerability indicator includes education, occupational position and a tenure component. The historical regression on modal assignments then reuses these inputs as regressors, sometimes with a different collapse. In particular, the 15.5 percentage point association for postsecondary education and assignment to the most digitally integrated profile is partly a comparison involving an input to the measurement model. It does not show that education predicts membership in an independently constructed class. Correcting classification error alone would not remove this overlap. Domestic use and remittance disposition also contribute to the historical use indicator, and reported transaction problems contribute to the adverse-experience indicator. Comparisons on these defining items help describe the construction; they cannot validate it externally \citep{bakkKuha2021}.
